\chapter{Geometric Multigrid: The V-Cycle}
\label{ch:multigrid}

\chapterorient{A smoother (\cref{ch:smoothers}) is a bargain that only half
works: a few cheap sweeps annihilate the jagged, high-frequency part of the error
and then stall, because what remains is \emph{smooth}---and smooth error is exactly
what a smoother cannot touch. This chapter turns that stall into a strategy.
The trick is a change of perspective: error that looks smooth on a fine grid looks
\emph{oscillatory} on a coarser one, where a smoother can kill it again, and far
more cheaply. Apply the idea once and you have a two-grid method; apply it to
itself---recursively, all the way down to a grid so small the problem is trivial---and
you have the V-cycle, the central algorithm of this chapter and one of the great
ideas of numerical computing. By the end you will be able to read the V-cycle as a
short recursive function, explain from the complementary-frequencies picture
\emph{why} it works, and state the property that makes it the optimal
preconditioner of \cref{ch:precond}: a convergence rate that does not degrade as the
mesh is refined, at a cost that grows only linearly in the number of unknowns.}

\section{The trouble a smoother leaves behind}

\Cref{ch:smoothers} studied the classical \emph{smoothers}---weighted Jacobi,
Gauss--Seidel, Chebyshev---the cheap stationary iterations that, applied to a linear
system $\mat{A}\vx=\vb$ arising from a discretized elliptic PDE, contract some
components of the error far faster than others. Recall the central fact established
there. Write the error after $k$ sweeps as $\ve_k = \vx - \vx_k$, where $\vx_k$ is
the current iterate and $\vx$ the exact solution. Expand $\ve_k$ in the eigenvectors
of the iteration: for a model problem each eigenvector is a \emph{Fourier mode}, a
sinusoid of some spatial frequency, and one sweep of a damped smoother multiplies
the mode of frequency $\theta$ by a factor $g(\theta)$ with $\abs{g(\theta)}<1$.

The defining property of a good smoother is that $\abs{g(\theta)}$ is small---near
zero---for the \emph{high} frequencies and close to $1$ for the \emph{low} ones.
A weighted-Jacobi sweep on the 1-D Laplacian, for instance, damps the most
oscillatory mode by a factor near $\tfrac13$ but the smoothest mode by a factor like
$1-O(h^2)$, agonizingly close to $1$. After a handful of sweeps the high-frequency
error is gone and the low-frequency error is essentially untouched. The iteration
appears to converge for three or four steps and then crawls.

\begin{keyidea}
A smoother is \emph{half} a solver. It cheaply annihilates the high-frequency
(oscillatory) part of the error and leaves the low-frequency (smooth) part almost
unchanged. The error that survives a few smoothing sweeps is, by construction,
\emph{smooth}. That residual smoothness is not a failure to be patched---it is a
\emph{signal}: smooth error is the part a coarser grid can represent, and on a
coarser grid it is no longer smooth.
\end{keyidea}

This is the opening every multigrid method exploits. The error left behind is smooth,
and smoothness is a \emph{grid-relative} notion. A function that wiggles slowly across
a fine mesh of $N$ points wiggles \emph{quickly} across a coarse mesh of $N/8$ points
covering the same domain---because there are fewer points for it to wiggle across. The
smoother that stalled on the fine grid will bite again on the coarse grid, where the
same error is now high-frequency. And the coarse grid is cheaper. The entire method is
the disciplined pursuit of this one observation.

\section{Complementary frequencies: the heart of the matter}
\label{sec:complementary}

Make the observation precise. Take the 1-D model: the domain $[0,1]$, a fine grid of
$n$ interior points at spacing $h = 1/(n+1)$, and a coarse grid of every \emph{other}
fine point, spacing $2h$. A Fourier mode on the fine grid is
\[
  \phi^h_j(x) = \sin(j\pi x), \qquad j = 1, 2, \dots, n,
\]
sampled at the fine nodes $x = h, 2h, \dots, nh$. The index $j$ is the
\emph{frequency}: $j=1$ is the smoothest half-sine; $j=n$ is the most oscillatory mode
the grid can resolve, alternating sign at every node. We split the modes at the
midpoint $n/2$:
\begin{itemize}[nosep]
\item \textbf{low (smooth) modes}, $1 \le j < n/2$: oscillate slowly; the smoother
  barely touches them ($\abs{g}\approx 1$);
\item \textbf{high (oscillatory) modes}, $n/2 \le j \le n$: oscillate rapidly; the
  smoother annihilates them ($\abs{g}\ll 1$).
\end{itemize}

Now ask the decisive question: what does a \emph{low} fine-grid mode look like when we
sample it on the coarse grid, at the nodes $x = 2h, 4h, \dots$ alone? The coarse grid
has half as many points, so it can only resolve frequencies up to $n/2$. A smooth
fine-grid mode, $j < n/2$, is \emph{also} a legitimate mode on the coarse grid---but
relative to the coarse grid's own resolution limit $n/2$, a mode at $j$ that was
``low'' on the fine grid sits a larger fraction of the way up toward the coarse limit.
The smoothest fine modes, the ones that defeated the smoother, become genuine,
killable, mid-to-high frequency content on the coarse grid. \Cref{fig:complementary}
draws exactly this: one smooth error sampled on both grids, oscillatory-looking on the
coarse one.

\begin{figure}[t]
\centering
\begin{tikzpicture}
\begin{axis}[
  width=12.8cm, height=5.4cm,
  xlabel={position $x$}, ylabel={error $e(x)$},
  xmin=0, xmax=1, ymin=-1.25, ymax=1.45,
  xtick={0,0.25,0.5,0.75,1}, ytick={-1,0,1},
  axis lines=left, tick align=outside,
  legend style={at={(0.5,1.02)}, anchor=south, legend columns=2,
    font=\footnotesize, draw=simGray!50},
  every axis plot/.append style={thick},
]
  % the underlying smooth error: a slowly-varying field (looks tame everywhere)
  \addplot[simBlue, very thick, samples=200, domain=0:1]
    {0.95*sin(deg(3.5*pi*x))*exp(-0.35*x)};
  \addlegendentry{smooth error (the continuous field)}
  % fine-grid samples: dense, follow the curve -> looks smooth
  \addplot[only marks, mark=*, mark size=1.0pt, simGreen]
    coordinates {
      (0.04,0.395)(0.08,0.685)(0.12,0.846)(0.16,0.860)(0.20,0.730)
      (0.24,0.479)(0.28,0.150)(0.32,-0.205)(0.36,-0.521)(0.40,-0.742)
      (0.44,-0.829)(0.48,-0.766)(0.52,-0.566)(0.56,-0.262)(0.60,0.087)
      (0.64,0.414)(0.68,0.661)(0.72,0.785)(0.76,0.764)(0.80,0.603)
      (0.84,0.334)(0.88,0.006)(0.92,-0.318)(0.96,-0.584)};
  \addlegendentry{fine-grid samples (smooth)}
  % coarse-grid samples: every other -> when connected, looks oscillatory
  \addplot[simRust, thick, mark=square*, mark size=1.7pt, densely dashed]
    coordinates {
      (0.04,0.395)(0.12,0.846)(0.20,0.730)(0.28,0.150)(0.36,-0.521)
      (0.44,-0.829)(0.52,-0.566)(0.60,0.087)(0.68,0.661)(0.76,0.764)
      (0.84,0.334)(0.92,-0.318)};
  \addlegendentry{coarse-grid samples (now oscillatory)}
\end{axis}
\end{tikzpicture}
\caption[Complementary frequencies]{The single idea behind multigrid. The blue curve
is a smooth error---the kind a smoother leaves behind, varying slowly across the
domain. Sampled on the \emph{fine} grid (green dots, every node), it traces the smooth
curve faithfully and a smoother sees a low-frequency mode it cannot damp. Sampled on
the \emph{coarse} grid (rust squares, every \emph{other} node, connected by the dashed
line), the \emph{same} error now swings sharply from sample to sample---it has become a
mid-to-high frequency mode relative to the coarse grid's coarser resolution, exactly
the kind a smoother \emph{can} annihilate. Smoothness is grid-relative; the coarse grid
turns the smoother's blind spot into its target.}
\label{fig:complementary}
\end{figure}

\begin{notationbox}
We will say a mode is \textbf{smooth} or \textbf{oscillatory} \emph{relative to a
grid}, never absolutely. A fixed error field is smooth on a fine grid and oscillatory
on a coarse one. ``High frequency'' means ``near the resolution limit of the grid in
question.'' The whole method rests on the fact that the smoother's reach---which modes
it damps---is tied to the grid's resolution limit, so changing the grid changes which
modes are in reach.
\end{notationbox}

This is the \emph{complementary-frequencies principle}. The smoother and the coarse
grid have complementary competences: the smoother handles everything oscillatory on the
fine grid; the coarse grid handles everything smooth on the fine grid, by re-presenting
it as something oscillatory that \emph{its} smoother can handle. Between them they cover
the whole spectrum. Neither alone is a solver; together they are. The two-grid method
of the next section is the minimal machine that makes them cooperate.

\section{The two-grid method}
\label{sec:twogrid}

Suppose we have a fine grid carrying the system $\mat{A}\vx=\vb$ and a single coarse
grid. We want a method that smooths on the fine grid (killing oscillatory error) and
\emph{also} corrects the smooth error by working on the coarse grid. The obstacle is
that the smooth error lives in the fine-grid solution, while the coarse grid speaks a
different, smaller vocabulary of vectors. We need a way to carry information between the
two---and an equation on the coarse grid for the coarse grid to solve.

\subsection{The residual equation}

The bridge is the \emph{residual equation}, which we have met before but which becomes
load-bearing here. Let $\vy$ be the current fine-grid approximation, after some
smoothing. Its error is $\ve = \vx - \vy$, unknown (it contains the exact solution). But
its \emph{residual},
\[
  \vr = \vb - \mat{A}\vy,
\]
is fully computable, and error and residual are linked by the exact linear identity
\[
  \mat{A}\,\ve = \mat{A}\vx - \mat{A}\vy = \vb - \mat{A}\vy = \vr.
\]
So $\mat{A}\,\ve = \vr$: the error solves the \emph{same} operator equation as the
original, but with the residual on the right. This is the \emph{error equation}. If we
could solve it we would have $\ve$ exactly and could finish in one step, $\vx = \vy +
\ve$. We cannot---solving $\mat{A}\,\ve=\vr$ on the fine grid is the original problem
again. But after smoothing, $\ve$ is \emph{smooth}, and a smooth $\ve$ is exactly what a
\emph{coarse} grid can represent well. So we solve the error equation
\emph{approximately, on the coarse grid}, and use the coarse solution to correct $\vy$.

\subsection{Restriction and prolongation}

To pose the error equation on the coarse grid we need two transfer operators
(\cref{fig:transfer}):
\begin{itemize}[nosep]
\item \textbf{Restriction} $\mat{R}\colon \text{fine}\to\text{coarse}$, which carries the
  fine-grid residual $\vr$ down to a coarse-grid residual $\vr_c = \mat{R}\vr$;
\item \textbf{Prolongation} $\mat{P}\colon \text{coarse}\to\text{fine}$, which carries a
  coarse-grid correction $\ve_c$ up to a fine-grid correction $\mat{P}\ve_c$.
\end{itemize}
Prolongation is \emph{interpolation}: given values at the coarse nodes, fill in the fine
nodes between them. In 1-D, linear interpolation copies the value at each coarse node to
the fine node sitting on it and averages the two neighbours for the fine node in between:
\[
  (\mat{P}\ve_c)_{2i} = (\ve_c)_i, \qquad
  (\mat{P}\ve_c)_{2i+1} = \tfrac12\bigl((\ve_c)_i + (\ve_c)_{i+1}\bigr).
\]
The matrix $\mat{P}$ is tall and thin: $N_f \times N_c$ with $N_f > N_c$, one column per
coarse node.

\begin{figure}[t]
\centering
\begin{tikzpicture}[scale=1.0, every node/.style={font=\scriptsize}]
  % fine grid (top), coarse grid (bottom)
  \def\fy{1.9} \def\cy{0}
  % fine nodes 0..8
  \foreach \i in {0,...,8}{
    \fill[simGreen] (\i*1.0,\fy) circle (2.0pt);}
  \draw[simGray] (0,\fy)--(8,\fy);
  \node[simGreen, anchor=south] at (4,\fy+0.18) {fine grid ($N_f$ nodes, spacing $h$)};
  % coarse nodes at 0,2,4,6,8
  \foreach \i in {0,2,4,6,8}{
    \fill[simRust] (\i*1.0,\cy) circle (2.6pt);}
  \draw[simGray] (0,\cy)--(8,\cy);
  \node[simRust, anchor=north] at (4,\cy-0.18) {coarse grid ($N_c$ nodes, spacing $2h$)};
  % prolongation (coarse -> fine), left half: interpolation arrows
  \begin{scope}
    % coincident copies
    \foreach \i in {0,2,4}{
      \draw[-{Stealth[length=1.6mm]}, simBlue, thick] (\i,\cy+0.12)--(\i,\fy-0.12);}
    % interpolated midpoints (averages of two neighbours)
    \foreach \i in {1,3}{
      \draw[-{Stealth[length=1.4mm]}, simBlue, opacity=0.55] (\i-1,\cy+0.12)--(\i,\fy-0.12);
      \draw[-{Stealth[length=1.4mm]}, simBlue, opacity=0.55] (\i+1,\cy+0.12)--(\i,\fy-0.12);}
    \node[simBlue, anchor=east, font=\footnotesize\bfseries] at (-0.25,1.0)
      {$\mat{P}$ (prolong)};
  \end{scope}
  % restriction (fine -> coarse), right half: weighted gather arrows
  \begin{scope}
    \foreach \i in {6,8}{
      \draw[-{Stealth[length=1.6mm]}, simRust, thick, densely dashed] (\i,\fy-0.12)--(\i,\cy+0.12);}
    \draw[-{Stealth[length=1.4mm]}, simRust, opacity=0.6, densely dashed] (5,\fy-0.12)--(6,\cy+0.12);
    \draw[-{Stealth[length=1.4mm]}, simRust, opacity=0.6, densely dashed] (7,\fy-0.12)--(6,\cy+0.12);
    \node[simRust, anchor=west, font=\footnotesize\bfseries] at (8.25,1.0)
      {$\mat{R}=\mat{P}^{\mathsf{T}}$ (restrict)};
  \end{scope}
\end{tikzpicture}
\caption[Prolongation and restriction]{The two grid-transfer operators between a fine
grid (green, spacing $h$) and a coarse grid (rust, spacing $2h$) holding every other
node. \emph{Left:} prolongation $\mat{P}$ (blue, solid arrows) is interpolation---each
coarse node's value is copied to the coincident fine node, and each in-between fine node
takes the average of its two coarse neighbours. \emph{Right:} restriction
$\mat{R}=\mat{P}^{\mathsf T}$ (rust, dashed arrows) is the transpose---it gathers each
coarse node's value as a weighted sum of nearby fine values, with weights mirroring the
interpolation stencil. Making restriction the transpose of prolongation is what keeps the
coarse operator symmetric and the whole method variationally consistent.}
\label{fig:transfer}
\end{figure}

Restriction is the reverse trip, and the natural choice---the one that makes the algebra
work out---is the \emph{transpose} of prolongation, up to scaling:
\[
  \mat{R} = \mat{P}^{\mathsf T}.
\]
Why the transpose? Because residuals and corrections are dual objects. A correction is a
\emph{primal} vector (it adds to the solution); a residual is a \emph{dual} vector (it
pairs with corrections through $\inner{\vr}{\ve}$ to measure work). Interpolating a
primal vector \emph{up} with $\mat{P}$ forces the dual vector to go \emph{down} with
$\mat{P}^{\mathsf T}$, so that the pairing $\inner{\vr}{\mat{P}\ve_c} =
\inner{\mat{P}^{\mathsf T}\vr}{\ve_c}$ is preserved between the grids. We will see in the
next subsection that this single choice is exactly what makes the coarse operator
symmetric.

\subsection{The Galerkin coarse operator}

The error equation on the fine grid is $\mat{A}\,\ve=\vr$. We want its coarse-grid
counterpart, $\mat{A}_c\,\ve_c = \vr_c$, with $\vr_c=\mat{P}^{\mathsf T}\vr$. What should
$\mat{A}_c$ be? Substitute the ansatz that the fine correction is the prolonged coarse
correction, $\ve \approx \mat{P}\ve_c$, into $\mat{A}\,\ve = \vr$ and restrict both
sides:
\[
  \mat{P}^{\mathsf T}\mat{A}\,\mat{P}\,\ve_c = \mat{P}^{\mathsf T}\vr.
\]
The operator on the left \emph{is} the right coarse operator:
\begin{equation}
  \mat{A}_c \;=\; \mat{P}^{\mathsf T}\mat{A}\,\mat{P}.
  \label{eq:galerkin}
\end{equation}
This is the \emph{Galerkin coarse operator}, or \emph{Galerkin coarsening}. It is not the
operator you would get by re-discretizing the PDE on the coarse grid (though for simple
problems the two agree up to scaling); it is the operator \emph{induced} by the fine
operator and the transfer maps. \Cref{eq:galerkin} has two virtues that no other choice
of $\mat{A}_c$ shares.

\begin{lemma}[Symmetry and definiteness are inherited]
\label{lem:galerkin}
If $\mat{A}$ is symmetric positive definite and $\mat{P}$ has full column rank, then the
Galerkin coarse operator $\mat{A}_c = \mat{P}^{\mathsf T}\mat{A}\mat{P}$ is symmetric
positive definite.
\end{lemma}
\begin{proof}
Symmetry: $\mat{A}_c^{\mathsf T} = (\mat{P}^{\mathsf T}\mat{A}\mat{P})^{\mathsf T} =
\mat{P}^{\mathsf T}\mat{A}^{\mathsf T}\mat{P} = \mat{P}^{\mathsf T}\mat{A}\mat{P} =
\mat{A}_c$, using $\mat{A}^{\mathsf T}=\mat{A}$. Definiteness: for any coarse vector
$\vv\ne\mathbf 0$, $\inner{\vv}{\mat{A}_c\vv} = \inner{\vv}{\mat{P}^{\mathsf
T}\mat{A}\mat{P}\vv} = \inner{\mat{P}\vv}{\mat{A}(\mat{P}\vv)} > 0$, because $\mat{P}\vv
\ne \mathbf 0$ (full column rank) and $\mat{A}$ is positive definite.
\end{proof}

So the coarse problem is again an SPD system: every tool we have for SPD systems---the
conjugate gradient method of \cref{ch:cg}, another round of smoothing, another level of
coarsening---applies to it unchanged. This self-similarity is what makes recursion
possible. The second virtue is \emph{variational consistency}: with $\mat{R} =
\mat{P}^{\mathsf T}$ and $\mat{A}_c=\mat{P}^{\mathsf T}\mat{A}\mat{P}$, the coarse-grid
correction is the best possible correction in the range of $\mat P$, measured in the
energy norm. The coarse grid does not merely \emph{approximate} the smooth-error
correction; it computes the \emph{optimal} one available to it.

\begin{physicsbox}
For finite elements the transfer operators are not invented---they are \emph{given by the
function-space hierarchy of \cref{ch:functionspaces}}. When the coarse FE space $V_H$ is a
genuine subspace of the fine space $V_h$ (a coarser mesh, or a lower polynomial order,
nested inside the finer one), every coarse basis function is exactly a linear combination
of fine basis functions. The coefficients of that combination \emph{are} the columns of
$\mat{P}$. Prolongation is then the natural inclusion $V_H \hookrightarrow V_h$, restriction
its transpose, and \cref{eq:galerkin} is automatic. Palace builds this nested hierarchy
once---an $h$-refined or $p$-coarsened tower of $\Hcurl$ (or $\Hone$) spaces---and reads
the prolongations straight off it. The next section makes that hierarchy the spine of the
recursion.
\end{physicsbox}

\subsection{The two-grid cycle assembled}

We now have every piece. The two-grid cycle, for one pass, is:
\begin{enumerate}[nosep]
\item \textbf{Presmooth:} apply a few smoother sweeps to $\mat{A}\vx=\vb$, starting from
  the current $\vx$, killing oscillatory error and leaving smooth error.
\item \textbf{Residual:} compute $\vr = \vb - \mat{A}\vx$.
\item \textbf{Restrict:} $\vr_c = \mat{P}^{\mathsf T}\vr$, carrying the residual to the
  coarse grid.
\item \textbf{Coarse solve:} solve $\mat{A}_c\,\ve_c = \vr_c$ on the coarse grid for the
  coarse correction.
\item \textbf{Prolong and correct:} $\vx \mathrel{+}= \mat{P}\,\ve_c$, adding the
  interpolated correction to the fine solution.
\item \textbf{Postsmooth:} a few more smoother sweeps, cleaning up the high-frequency
  error that interpolation reintroduced.
\end{enumerate}
Steps 1 and 6 handle the oscillatory error; steps 2--5 handle the smooth error via the
coarse grid. This is a complete, correct method---and \cref{wex:twogrid} below runs it on
real numbers and watches the error collapse. But it has one unsatisfying feature: step 4
says ``solve the coarse problem,'' and for a large fine grid the coarse problem is still
large. The resolution of that complaint is the idea that makes multigrid \emph{multi}.

\section{The V-cycle: recursion down the grids}
\label{sec:vcycle}

Step 4 of the two-grid cycle asks us to solve $\mat{A}_c\,\ve_c=\vr_c$ on the coarse grid.
But $\mat{A}_c$ is SPD (\cref{lem:galerkin}), so it is a problem of exactly the same kind
as the original---and we know a good way to attack a problem of that kind: the two-grid
cycle itself. Rather than solve the coarse problem exactly, \emph{apply one two-grid cycle
to it}, using a still-coarser grid below it. That two-grid cycle in turn defers its own
coarse solve to a yet-coarser grid, and so on, until we reach a grid so small---a handful
of unknowns---that solving directly is trivial. The two-grid method, applied to itself all
the way down, is the \textbf{V-cycle}.

\begin{keyidea}
\textbf{Smooth on the fine grid; solve the rest on a coarser grid---recursively.}
On every level, presmooth to kill that level's oscillatory error, push the smooth residual
down to the next-coarser level, let that level handle it the same way, bring the correction
back up, and postsmooth. The recursion bottoms out on a grid small enough to solve
directly. Because each level only ever does cheap smoothing (the expensive ``solve'' is
always deferred downward), the total cost is a geometric series that sums to $O(N)$, and
because each level kills exactly the frequencies it is responsible for, the convergence
rate is \emph{independent of the mesh size}. That conjunction---$O(N)$ work and
$h$-independent convergence---is what the word \emph{optimal} means in \cref{ch:precond}.
\end{keyidea}

\subsection{The V-cycle as a recursive function}

Number the levels $0, 1, \dots, L$, with $L$ the finest (the grid we actually want to
solve on) and $0$ the coarsest (the trivial base case). Each level $\ell$ carries its own
operator $\mat{A}_\ell$, its own smoother $\mat{B}_\ell$, and---for $\ell>0$---a
prolongation $\mat{P}_{\ell-1}$ mapping from level $\ell-1$ up to level $\ell$. The
V-cycle is then a single recursive function of the level $\ell$ and the current iterate
$\vx$. \Cref{lst:vcycle} writes it as a pure function---no in-place mutation, every step a
tensor-in/tensor-out map---in the calculus of \cref{part:framework}; this is the form Palace's
\code{GeometricMultigridSolver} takes once its scratch-vector mutations are rotated away.

\begin{lstlisting}[style=l4style, caption={The V-cycle as a level-recursive pure function. Each call presmooths, restricts the residual, recurses on the coarser level, prolongs the returned correction, and postsmooths. The base case $\ell=0$ solves directly.}, label={lst:vcycle}]
-- one V-cycle pass at level l: improve x toward solving A[l] x = b
vcycle :: [LinOp] -> [Smoother] -> [LinOp] -> Solver
       -> Int -> Tensor[N] -> Tensor[N] -> Tensor[N]
vcycle as bs ps coarseSolve 0 b x =
    coarseSolve b                          -- base case: solve A[0] e = b exactly
vcycle as bs ps coarseSolve l b x =
  let y   = presmooth (bs!l) b x           -- 1. kill oscillatory error on level l
      r   = b `sub` apply (as!l) y         -- 2. residual r = b - A[l] y
      rc  = apply_transpose (ps!(l-1)) r   -- 3. restrict  r_c = P^T r   (fine -> coarse)
      ec  = vcycle as bs ps coarseSolve (l-1) rc (zeros)   -- 4. RECURSE: solve A[l-1] e_c = r_c
      y'  = y `add` apply (ps!(l-1)) ec    -- 5. prolong & correct  y += P e_c
  in  postsmooth (bs!l) b y'               -- 6. clean up interpolation error
\end{lstlisting}

Read the function once more with the two-grid cycle of \cref{sec:twogrid} beside it: the
six numbered comments are the six numbered steps, with step~4---formerly ``solve the coarse
problem''---now a \emph{recursive call to \code{vcycle} itself} one level down. The base
case \code{vcycle ... 0 b x} does the only genuine solve in the whole algorithm, and it is
cheap because level $0$ is tiny. Everything above the base case does nothing but smooth,
restrict, and prolong. The coarse operators $\mat{A}_\ell$ and prolongations
$\mat{P}_{\ell-1}$ are precomputed once from the function-space hierarchy
(\cref{ch:functionspaces}); the recursion just walks them.

\begin{notationbox}
Each leg of the V-cycle---presmooth, the coarse-grid correction, postsmooth---has the same
algebraic shape: take the current guess $\vy$, compute its residual $\vr=\vb-\mat{A}\vy$,
apply some approximate inverse $\mat{B}$ to the residual, and add the result back:
\[
  \vy \;\longmapsto\; \vy + \mat{B}\,(\vb - \mat{A}\vy).
\]
This is the \emph{correction step} of \cref{ch:precond}: a presmooth leg takes $\mat{B}$ to
be the point smoother; the coarse-grid-correction leg takes $\mat{B} =
\mat{P}\,\mat{A}_c^{-1}\mat{P}^{\mathsf T}$ (smooth the residual by solving on the coarse
grid). The V-cycle is a \emph{composition of correction steps}, alternating point smoothers
with one coarse-grid smoother---which is itself a composition of correction steps one level
down. Reading multigrid this way---one combinator, applied at every level and recursively
inside itself---is the calculus view of \cref{part:framework}.
\end{notationbox}

\subsection{Why it is called a V}

Trace the flow of control through one call to \code{vcycle} at the finest level $L$.
Presmoothing happens on the way \emph{down}: level $L$ smooths, then hands a residual to
level $L-1$, which smooths and hands a residual to $L-2$, and so on down to level $0$.
That is the descending stroke of the letter. At level $0$ the direct solve happens---the
point of the V. Then the corrections climb back \emph{up}: level $0$ returns a correction to
$1$, which prolongs it, postsmooths, and returns to $2$, up to the finest level. That is the
ascending stroke. Drawn with the levels stacked vertically and time running left to right,
the path of the computation is a literal V (\cref{fig:vcycle}). This is the hero diagram of
the chapter; every multigrid method is a variation on its shape.

\begin{figure}[p]
\centering
\begin{tikzpicture}[scale=1.0, every node/.style={font=\footnotesize}]
  % vertical: level (fine at top). horizontal: time.
  % level y-coordinates
  \def\Lf{4.5} \def\La{3.0} \def\Lb{1.5} \def\Lc{0.0}
  % level guide lines + labels
  \foreach \y/\lab/\dof in {%
    4.5/{level $L$ (finest)}/{$N$ dof},%
    3.0/{level $L-1$}/{$N/2^d$},%
    1.5/{level $1$}/{$\vdots$},%
    0.0/{level $0$ (coarsest)}/{tiny}%
  }{
    \draw[simGray!35, densely dotted] (0.2,\y)--(11.3,\y);
    \node[anchor=east, simGray, font=\scriptsize] at (0.1,\y) {\lab};
    \node[anchor=west, simGray!80, font=\scriptsize\itshape] at (11.4,\y) {\dof};
  }
  % the V path: descending (presmooth+restrict), then ascending (prolong+postsmooth)
  % node coordinates along the V
  \coordinate (d0) at (1.6,\Lf);   % presmooth finest
  \coordinate (d1) at (3.4,\La);   % presmooth L-1
  \coordinate (d2) at (5.2,\Lb);   % presmooth 1
  \coordinate (bot) at (6.4,\Lc);  % coarse solve
  \coordinate (u2) at (7.6,\Lb);   % postsmooth 1
  \coordinate (u1) at (9.4,\La);   % postsmooth L-1
  \coordinate (u0) at (11.2,\Lf);  % postsmooth finest
  % descending strokes (restriction)
  \draw[-{Stealth[length=2.4mm]}, simBlue, very thick] (d0)--(d1);
  \draw[-{Stealth[length=2.4mm]}, simBlue, very thick] (d1)--(d2);
  \draw[-{Stealth[length=2.4mm]}, simBlue, very thick] (d2)--(bot);
  % ascending strokes (prolongation)
  \draw[-{Stealth[length=2.4mm]}, simRust, very thick] (bot)--(u2);
  \draw[-{Stealth[length=2.4mm]}, simRust, very thick] (u2)--(u1);
  \draw[-{Stealth[length=2.4mm]}, simRust, very thick] (u1)--(u0);
  % node markers
  \foreach \pt/\col in {d0/simGreen, d1/simGreen, d2/simGreen,
                        u2/simGold, u1/simGold, u0/simGold}{
    \fill[\col] (\pt) circle (3pt);}
  \fill[simViolet] (bot) circle (4pt);
  % leg labels (descending = presmooth + restrict)
  \node[simGreen, anchor=south, font=\scriptsize] at (d0) {presmooth};
  \node[simBlue, anchor=south east, font=\scriptsize] at ($(d0)!0.5!(d1)$) {restrict $\mat{P}^{\mathsf T}\vr$};
  \node[simBlue, anchor=south east, font=\scriptsize] at ($(d1)!0.5!(d2)$) {restrict};
  \node[simBlue, anchor=east, font=\scriptsize] at ($(d2)!0.5!(bot)+(0.1,0)$) {restrict};
  % bottom
  \node[simViolet, anchor=north, font=\scriptsize\bfseries] at (bot) {coarse solve};
  % ascending labels (prolong + postsmooth)
  \node[simRust, anchor=west, font=\scriptsize] at ($(bot)!0.5!(u2)-(0.1,0)$) {prolong $\mat{P}\ve_c$};
  \node[simRust, anchor=south west, font=\scriptsize] at ($(u2)!0.5!(u1)$) {prolong};
  \node[simRust, anchor=south west, font=\scriptsize] at ($(u1)!0.5!(u0)$) {prolong};
  \node[simGold, anchor=south, font=\scriptsize] at (u0) {postsmooth};
  % the big V annotation
  \node[simInk, font=\large\bfseries, opacity=0.18] at (6.4,2.5) {V};
\end{tikzpicture}
\caption[The V-cycle]{The hero diagram. Levels stacked vertically (finest at top,
coarsest at bottom), time running left to right. The computation descends the left stroke
of a V---\emph{presmoothing} (green) on each level, then \emph{restricting} the residual
(blue arrows, $\mat{P}^{\mathsf T}\vr$) to the next-coarser level---until it reaches the
bottom, where the coarsest problem is solved directly (violet). It then climbs the right
stroke---\emph{prolonging} each correction (rust arrows, $\mat{P}\ve_c$) up to the next-finer
level and \emph{postsmoothing} (gold)---back to the finest level. Each level does only cheap
smoothing and grid transfers; the single genuine solve happens on the tiny coarsest grid.
The grid shrinks by a factor $2^d$ per level (right margin), so the work per level shrinks
geometrically and the total is $O(N)$.}
\label{fig:vcycle}
\end{figure}

\subsection{Building the level hierarchy}

The V-cycle presupposes a tower of grids and the operators connecting them. Where does the
tower come from? \Cref{ch:functionspaces} already built it. Palace constructs a
\emph{finite-element space hierarchy}: a coarse-to-fine stack of nested FE spaces,
\[
  V_0 \subset V_1 \subset \cdots \subset V_L,
\]
obtained either by $h$-refinement (each $V_{\ell}$ a finer mesh than $V_{\ell-1}$) or by
$p$-coarsening (each $V_{\ell}$ a higher polynomial order). Because the spaces are nested,
each inclusion $V_{\ell-1}\hookrightarrow V_\ell$ \emph{is} a prolongation
$\mat{P}_{\ell-1}$, read directly off the hierarchy. The hierarchy is built once, before the
solve, by folding a level-adding operation over the refinement sequence---a coarse seed
space, then one level per refinement step---and thereafter the V-cycle consumes it
read-only. \Cref{fig:hierarchy} shows the geometric picture: the same domain meshed at three
resolutions, each coarse mesh nested in the finer one. The smoothers $\mat{B}_\ell$ and the
Galerkin operators $\mat{A}_\ell = \mat{P}_{\ell-1}^{\mathsf T}\mat{A}_{\ell+1}\mat{P}_{\ell-1}$
are likewise precomputed per level. With the tower in hand, the recursion of
\cref{lst:vcycle} is all that remains.

\begin{figure}[t]
\centering
\begin{tikzpicture}[scale=1.0, every node/.style={font=\scriptsize}]
  % three nested meshes of the same square domain
  \foreach \px/\n/\ttl/\col in {0/4/{fine ($h$)}/simGreen,
                                 5/2/{medium ($2h$)}/simBlue,
                                 10/1/{coarse ($4h$)}/simRust}{
    \begin{scope}[shift={(\px,0)}]
      \draw[simInk, thick] (0,0) rectangle (3,3);
      % grid lines
      \pgfmathsetmacro{\step}{3/\n}
      \foreach \k in {1,...,\n}{
        \pgfmathsetmacro{\pos}{\k*\step}
        \ifnum\k<\n
          \draw[\col!70] (\pos,0)--(\pos,3);
          \draw[\col!70] (0,\pos)--(3,\pos);
        \fi
      }
      % nodes
      \foreach \i in {0,...,\n}{\foreach \j in {0,...,\n}{
        \fill[\col] (\i*\step,\j*\step) circle (1.6pt);}}
      \node[font=\footnotesize\bfseries, \col] at (1.5,-0.45) {\ttl};
    \end{scope}}
  % refinement arrows between
  \draw[-{Stealth[length=2.4mm]}, simViolet, thick] (8.7,1.5)--(8.2,1.5)
    node[midway, above, simViolet]{$\mat{P}$};
  \draw[-{Stealth[length=2.4mm]}, simViolet, thick] (3.7,1.5)--(3.2,1.5)
    node[midway, above, simViolet]{$\mat{P}$};
  \node[anchor=north, simGray, font=\scriptsize\itshape] at (6.5,-1.0)
    {prolongation interpolates from each grid to the next finer; restriction $\mat{P}^{\mathsf T}$ runs the other way};
\end{tikzpicture}
\caption[The grid hierarchy]{The multigrid hierarchy of nested grids on one square domain.
\emph{Right to left:} a coarse mesh (rust, spacing $4h$), a medium mesh (blue, $2h$), a fine
mesh (green, $h$). Each coarse grid's nodes are a subset of the finer grid's, so the coarse
finite-element space is nested in the finer one and the inclusion is the prolongation
$\mat{P}$ (violet arrows). For finite elements the nesting comes free from the
function-space hierarchy of \cref{ch:functionspaces}; the V-cycle of \cref{fig:vcycle}
walks up and down this tower.}
\label{fig:hierarchy}
\end{figure}

\section{Why it is optimal: $O(N)$ work, mesh-independent convergence}
\label{sec:optimal}

The V-cycle's reputation rests on two quantitative claims. Each is worth establishing
carefully, because together they are what \cref{ch:precond} means by an \emph{optimal}
preconditioner and what no other general-purpose solver delivers for elliptic problems.

\subsection{Linear work per cycle}

\emph{Claim: one V-cycle costs $O(N)$ operations, where $N$ is the number of fine-grid
unknowns.} The work on each level is proportional to the number of unknowns on that level:
a fixed number of smoother sweeps, one residual, one restriction, one prolongation, all of
which touch each unknown a constant number of times. Let $N_\ell$ be the unknown count on
level $\ell$, with $N_L = N$ the finest. In $d$ space dimensions, coarsening by a factor of
two per direction divides the unknown count by $2^d$ each level:
\[
  N_{L-1} \approx \frac{N}{2^d}, \quad N_{L-2}\approx \frac{N}{4^d}, \quad \dots
\]
The total work is the sum over levels, $W = c\sum_{\ell} N_\ell$, a geometric series
\[
  W \;\approx\; c\,N\Bigl(1 + 2^{-d} + 2^{-2d} + \cdots\Bigr)
    \;=\; c\,N\,\frac{1}{1 - 2^{-d}}.
\]
The ratio $1/(1-2^{-d})$ is a constant: $2$ in 1-D, $\tfrac43$ in 2-D, $\tfrac87$ in 3-D.
So the whole tower costs only a small constant factor more than the finest level alone---the
coarse levels are essentially free. \Cref{fig:costseries} draws the geometric series of
per-level costs; \cref{wex:cost} makes the sum concrete. \emph{One V-cycle is $O(N)$.}

\begin{figure}[t]
\centering
\begin{tikzpicture}
\begin{axis}[
  width=11.5cm, height=5.2cm,
  ybar, bar width=15pt,
  xlabel={multigrid level (finest $=L$, then coarser)},
  ylabel={work $\propto N_\ell$},
  symbolic x coords={L, L-1, L-2, L-3, L-4},
  xtick=data, ymin=0, ymax=1.15,
  ytick={0,0.25,0.5,0.75,1.0},
  yticklabel={\pgfmathprintnumber{\tick}\,$N$},
  axis lines=left, tick align=outside,
  nodes near coords, nodes near coords style={font=\scriptsize, simGray},
  every node near coord/.append style={/pgf/number format/fixed,
    /pgf/number format/precision=3},
  bar shift=0pt,
]
  % work per level in 3D: N, N/8, N/64, ...
  \addplot[fill=simBlue!70, draw=simBlue] coordinates {
    (L,1.0) (L-1,0.125) (L-2,0.0156) (L-3,0.00195) (L-4,0.00024)};
\end{axis}
\end{tikzpicture}
\caption[The per-level cost geometric series]{Work per multigrid level in 3-D ($d=3$),
in units of the finest-level work $N$. Each coarser level has $1/2^d = 1/8$ the unknowns of
the one above, so its cost is one-eighth as large: $N,\ N/8,\ N/64,\dots$. The bars shrink so
fast that the entire tower below the finest level adds only $1/(2^3-1) = 1/7$ to the finest
level's cost---the coarse grids are nearly free. Summed, the geometric series gives total
work $\tfrac{8}{7}N = O(N)$ per V-cycle, independent of how many levels the tower has.}
\label{fig:costseries}
\end{figure}

\subsection{Mesh-independent convergence}

\emph{Claim: the V-cycle reduces the error by a fixed factor $\rho<1$ per cycle, and
$\rho$ does not grow as the mesh is refined.} This is the deeper claim, and the
complementary-frequencies principle of \cref{sec:complementary} is its proof sketch. Split
the error into smooth and oscillatory parts relative to the fine grid. The smoother
contracts the oscillatory part by a factor bounded below $1$ \emph{independent of $h$}---a
property called \emph{smoothing}. The coarse-grid correction contracts the smooth part,
and by variational consistency (\cref{lem:galerkin}) it does so optimally, again with a
factor bounded below $1$ \emph{independent of $h$}---a property called \emph{approximation}.
Every error component is hit by one mechanism or the other with an $h$-independent factor,
so the overall contraction $\rho$ is bounded below $1$ uniformly in $h$. A typical
well-tuned V-cycle has $\rho \approx 0.1$ to $0.2$: each cycle kills an order of magnitude
of error, regardless of mesh size.

Contrast this with the unpreconditioned conjugate gradient method of \cref{ch:cg}. Its
iteration count to a fixed tolerance scales as $\sqrt{\kappa}$, where $\kappa$ is the
condition number of $\mat{A}$. For a second-order elliptic operator, $\kappa = O(h^{-2})$,
so $\sqrt{\kappa} = O(h^{-1})$: \emph{halving the mesh spacing doubles the iteration count.}
Multigrid's iteration count, by contrast, is \emph{flat}---a constant number of V-cycles
reaches tolerance no matter how fine the mesh. \Cref{fig:meshindep} draws the two curves;
\cref{wex:cost} puts numbers on them.

\begin{figure}[t]
\centering
\begin{tikzpicture}
\begin{axis}[
  width=11.5cm, height=5.6cm,
  xlabel={problem size $N$ (degrees of freedom)},
  ylabel={iterations to tolerance},
  xmode=log, log basis x=2,
  xmin=700, xmax=1.2e6, ymin=0, ymax=320,
  xtick={1024,4096,16384,65536,262144,1048576},
  xticklabels={$2^{10}$,$2^{12}$,$2^{14}$,$2^{16}$,$2^{18}$,$2^{20}$},
  ytick={0,50,100,150,200,250,300},
  axis lines=left, tick align=outside, grid=major, grid style={simGray!18},
  legend style={at={(0.03,0.97)}, anchor=north west, font=\footnotesize,
    draw=simGray!50},
  every axis plot/.append style={very thick},
]
  % unpreconditioned CG: iterations ~ sqrt(kappa) ~ 1/h ~ N^{1/d} (d=2 => N^{1/2})
  \addplot[simRust, mark=*, mark size=1.6pt] coordinates {
    (1024,9.6)(4096,19.2)(16384,38.4)(65536,76.8)(262144,153.6)(1048576,307.2)};
  \addlegendentry{unpreconditioned CG: $\sim\!\sqrt{\kappa}\sim 1/h$}
  % multigrid V-cycle: flat
  \addplot[simBlue, mark=square*, mark size=1.6pt] coordinates {
    (1024,8)(4096,8)(16384,8)(65536,9)(262144,9)(1048576,9)};
  \addlegendentry{multigrid V-cycle: flat ($h$-independent)}
\end{axis}
\end{tikzpicture}
\caption[Mesh-independent convergence]{The property that makes multigrid optimal.
\emph{Rust:} unpreconditioned conjugate gradient (\cref{ch:cg}) needs $\sim\sqrt{\kappa}\sim
1/h$ iterations, which \emph{doubles} every time the mesh is halved (note the log scale on
$N$)---refining the mesh makes the solver slower and slower. \emph{Blue:} the multigrid
V-cycle takes a nearly constant number of cycles regardless of mesh size: its convergence
factor $\rho<1$ is bounded \emph{independent} of $h$. This flatness is the definition of an
\emph{optimal} solver in \cref{ch:precond}, and multigrid is the canonical example.}
\label{fig:meshindep}
\end{figure}

\begin{keyidea}
The two claims combine into the headline. \emph{One V-cycle costs $O(N)$, and a fixed
number of V-cycles---independent of $N$---reaches any tolerance.} Therefore solving an
elliptic PDE to fixed accuracy costs $O(N)$ total: optimal, in the strict sense that you
cannot beat linear in the number of unknowns, since you must at least look at each unknown
once. No purely algebraic solver---not Gaussian elimination ($O(N^3)$, or $O(N^{1.5})$ for
2-D sparse), not unpreconditioned CG ($O(N^{1.5})$ in 2-D)---matches this. Multigrid is the
gold standard against which elliptic solvers are measured.
\end{keyidea}

\section{Worked examples}

\subsection{A two-grid cycle on 1-D Poisson}

We run the two-grid method of \cref{sec:twogrid} on a small concrete problem and watch the
error collapse. The arithmetic is deliberately tractable so every step can be checked by
hand.

\begin{workedexample}{A two-grid cycle on the 1-D Poisson equation}{twogrid}
Discretize $-u'' = f$ on $[0,1]$ with homogeneous Dirichlet conditions, on a fine grid of
$n=7$ interior points at spacing $h=\tfrac18$. The operator is the standard tridiagonal
$\mat{A} = h^{-2}\,\mathrm{tridiag}(-1, 2, -1)$, size $7\times 7$. Take a right-hand side
whose exact solution is the smooth function $u(x)=\sin(\pi x)$ (so the discrete error we
introduce will be smooth-plus-oscillatory and we can watch each part go).

\textbf{Initial error.} Start from an initial guess corrupted by both a smooth and an
oscillatory component: let the initial error be
$\ve_0 = \underbrace{\sin(\pi x)}_{\text{smooth, } j=1} +
\tfrac12\underbrace{\sin(7\pi x)}_{\text{oscillatory, } j=7}$,
sampled at the seven fine nodes. Its energy norm (we track $\norm{\ve}_\infty$ for
legibility) starts near $\norm{\ve_0}_\infty \approx 1.42$.

\textbf{Step 1 --- presmooth.} Apply two sweeps of weighted Jacobi, $\vx \leftarrow \vx +
\omega\mat{D}^{-1}(\vb-\mat{A}\vx)$ with $\omega=\tfrac23$ and $\mat{D}=2h^{-2}\mathrm{I}$.
The Jacobi smoothing factor on this operator is $|1-\tfrac23(1-\cos\theta_j)|$, which is
$\approx 0.95$ for the smooth mode $j=1$ but $\approx 0.10$ for the oscillatory mode $j=7$.
After two sweeps the oscillatory amplitude $\tfrac12$ has fallen to
$\tfrac12(0.10)^2 \approx 0.005$---gone---while the smooth amplitude $1$ has barely
moved, to $\approx 0.90$. The error is now $\norm{\ve}_\infty \approx 0.90$, and---this is
the point---it is now essentially \emph{pure smooth mode}.

\textbf{Step 2 --- residual.} Compute $\vr = \vb - \mat{A}\vx$ on the fine grid. Because the
surviving error is the smooth mode $\ve\approx 0.90\sin(\pi x)$ and $\mat{A}\sin(\pi x) =
\lambda_1\sin(\pi x)$ with $\lambda_1 = h^{-2}\cdot 2(1-\cos\pi h)\approx 9.74$, the residual
is itself smooth, $\vr = \mat{A}\ve \approx 8.8\sin(\pi x)$.

\textbf{Step 3 --- restrict.} The coarse grid has $n_c=3$ interior points at spacing
$2h=\tfrac14$. Full-weighting restriction $\mat{P}^{\mathsf T}$ (stencil
$\tfrac14[1,2,1]$) carries $\vr$ to a coarse residual $\vr_c$ of length $3$, still a clean
smooth coarse mode.

\textbf{Step 4 --- coarse solve.} Solve the $3\times 3$ system
$\mat{A}_c\,\ve_c = \vr_c$ \emph{exactly} (three unknowns---a trivial direct solve), where
$\mat{A}_c = \mat{P}^{\mathsf T}\mat{A}\mat{P}$ is the Galerkin operator,
itself the $3\times3$ tridiagonal Laplacian at spacing $2h$. The coarse correction $\ve_c$
recovers almost exactly the coarse samples of the smooth error $0.90\sin(\pi x)$.

\textbf{Step 5 --- prolong and correct.} Interpolate: $\vx \leftarrow \vx +
\mat{P}\,\ve_c$. The prolonged correction $\mat{P}\ve_c$ matches the smooth error $\ve$ to
within the interpolation error of a smooth function on a grid that resolves it well---a few
percent. After the correction the fine-grid error has dropped from $\approx 0.90$ to
$\approx 0.06$.

\textbf{Step 6 --- postsmooth.} Two more Jacobi sweeps clean up the small oscillatory error
that linear interpolation injected, bringing $\norm{\ve}_\infty \approx 0.02$.

\textbf{The tally.} One two-grid cycle drove the error from $1.42$ to $0.02$---a factor of
about $70$, two orders of magnitude---at the cost of a few fine-grid Jacobi sweeps plus one
trivial $3\times3$ solve. Smoothing alone, run for the same number of sweeps, would have
removed the oscillatory mode and then stalled at $\approx 0.90$, having made essentially no
progress on the smooth mode. The coarse grid did in one $3\times 3$ solve what no amount of
fine-grid smoothing could do. \emph{That} is the complementary-frequencies principle paying
off in numbers.
\end{workedexample}

\subsection{The work estimate and the iteration-count contrast}

The second example makes the two optimality claims of \cref{sec:optimal} quantitative: the
$O(N)$ per-cycle cost as a summed geometric series, and the mesh-independent iteration count
set against CG's growth.

\begin{workedexample}{$O(N)$ work and the CG contrast}{cost}
\textbf{Part (a): one V-cycle is $O(N)$.} Take a 3-D problem, $d=3$, with $N$ unknowns on the
finest grid. Coarsening halves the grid spacing's inverse---divides the unknown count by
$2^d = 8$---per level. Let the work on a level be $c$ operations per unknown (a fixed number
of smoother sweeps plus the transfers). The total work is
\[
  W = c\Bigl(N + \tfrac{N}{8} + \tfrac{N}{64} + \tfrac{N}{512} + \cdots\Bigr)
    = cN\sum_{k=0}^{\infty} 8^{-k} = cN\cdot\frac{1}{1-\tfrac18} = cN\cdot\frac{8}{7}.
\]
The sum is a finite constant, $\tfrac87 \approx 1.143$: \emph{the entire tower of coarse
grids adds only $14\%$ to the cost of the finest level alone.} So $W = O(N)$, with a small
constant. In 2-D the factor is $\tfrac{1}{1-1/4} = \tfrac43$; in 1-D, $\tfrac{1}{1-1/2}=2$.
In every dimension the geometric series converges and the per-cycle cost is strictly linear
in $N$.

\textbf{Part (b): total cost to solve.} If a fixed number $m$ of V-cycles---say $m\approx
8$---reaches the tolerance \emph{regardless of $N$} (the mesh-independence claim), then the
total solve cost is $m\cdot W = 8\cdot\tfrac87 cN = O(N)$. Solving the PDE costs linear time.

\textbf{Part (c): the CG contrast.} Now estimate unpreconditioned CG on the same family of
2-D problems. Its iteration count to fixed tolerance is $\approx \tfrac12\sqrt{\kappa}$ with
$\kappa \approx (h^{-2})\cdot(\text{const})$, so iterations $\propto 1/h \propto
\sqrt{N}$ in 2-D. Tabulate both as the mesh is refined, doubling resolution each row (so $N$
quadruples in 2-D):
\begin{center}\small
\begin{tabular}{@{}rrr@{}}
\toprule
$N$ & CG iterations $\sim\!\sqrt{N}$ & V-cycles \\
\midrule
$1\,024$      & $\approx 16$  & $8$ \\
$4\,096$      & $\approx 32$  & $8$ \\
$16\,384$     & $\approx 64$  & $8$ \\
$65\,536$     & $\approx 128$ & $9$ \\
$262\,144$    & $\approx 256$ & $9$ \\
\bottomrule
\end{tabular}
\end{center}
CG's iteration count \emph{doubles every time the mesh is halved}; multigrid's is flat. And
each CG iteration also costs $O(N)$ (one sparse mat-vec), so CG's total solve cost in 2-D is
$O(N\sqrt{N}) = O(N^{1.5})$---superlinear---against multigrid's $O(N)$. At $N=262\,144$
multigrid does the work of $9$ cycles while CG grinds through $256$ iterations; the gap
widens without bound as the mesh refines. This table \emph{is} \cref{fig:meshindep} in
numbers, and it is the entire case for preconditioning the Krylov solve with a V-cycle---the
subject of the next section.
\end{workedexample}

\section{Variants and deployment}
\label{sec:variants}

The V-cycle is the prototype, not the only shape. A few variations matter in practice, and
one deployment pattern---multigrid \emph{inside} a Krylov solver---is how Palace actually
uses it.

\subsection{V-cycle, W-cycle, and full multigrid}

The recursion of \cref{lst:vcycle} makes exactly \emph{one} recursive call per level. If the
coarse-grid correction is not accurate enough---because the coarse problem itself was solved
too crudely---one can make \emph{two} recursive calls per level instead, visiting each coarse
level twice. The resulting flow traces a \emph{W} rather than a V (\cref{fig:cycleshapes},
middle): more coarse-grid work, a stronger correction, useful for harder problems where a
single coarse visit leaves too much error. \emph{Full multigrid} (FMG) is a different idea
again: rather than start on the finest grid with a poor guess, start on the \emph{coarsest}
grid, solve it, prolong the solution to the next level as an excellent initial guess, do a
V-cycle there, prolong again, and so on up to the finest grid (\cref{fig:cycleshapes},
right). FMG reaches discretization accuracy in a single sweep up the tower---often the most
efficient option of all, because every level starts from a guess the level below already
made good.

\begin{figure}[t]
\centering
\begin{tikzpicture}[scale=0.82, every node/.style={font=\scriptsize}]
  % shared: 4 levels, y = 0..3 (0 coarsest)
  % --- V-cycle ---
  \begin{scope}
    \foreach \y in {0,1,2,3}{\draw[simGray!30, densely dotted] (0,\y)--(3.2,\y);}
    \node[font=\footnotesize\bfseries, simInk] at (1.6,3.7) {V-cycle};
    \draw[-{Stealth[length=2mm]}, simBlue, thick]
      (0.2,3)--(1.0,2)--(1.8,1)--(2.2,0)--(2.6,1)--(3.0,2)--(3.4,3);
    \draw[simBlue, thick]
      (0.2,3)--(1.0,2)--(1.8,1)--(2.2,0)--(2.6,1)--(3.0,2)--(3.4,3);
    \node[anchor=east, simGray, font=\scriptsize] at (-0.1,3) {fine};
    \node[anchor=east, simGray, font=\scriptsize] at (-0.1,0) {coarse};
  \end{scope}
  % --- W-cycle ---
  \begin{scope}[shift={(4.6,0)}]
    \foreach \y in {0,1,2,3}{\draw[simGray!30, densely dotted] (0,\y)--(4.4,\y);}
    \node[font=\footnotesize\bfseries, simInk] at (2.2,3.7) {W-cycle};
    \draw[simRust, thick]
      (0.2,3)--(0.9,2)--(1.4,1)--(1.7,0)--(2.0,1)--(2.3,0)--(2.7,1)--(3.2,2)
      --(3.6,1)--(3.9,0)--(4.2,1)--(4.6,3);
  \end{scope}
  % --- FMG ---
  \begin{scope}[shift={(10.0,0)}]
    \foreach \y in {0,1,2,3}{\draw[simGray!30, densely dotted] (0,\y)--(4.4,\y);}
    \node[font=\footnotesize\bfseries, simInk] at (2.2,3.7) {full multigrid};
    % start at coarsest, climb with V-cycles
    \draw[simGreen, thick]
      (0.2,0)--(0.6,1)--(0.4,0)            % V at level 1
      (0.6,1)--(1.2,2)--(0.9,1)--(0.7,0)--(1.0,1)--(1.4,2)   % V at level 2
      (1.4,2)--(2.2,3)--(1.8,2)--(1.5,1)--(1.3,0)--(1.6,1)--(2.0,2)--(2.6,3); % V at finest
    \node[anchor=north, simGray, font=\scriptsize\itshape] at (2.2,-0.5)
      {start coarse, climb};
  \end{scope}
\end{tikzpicture}
\caption[V-cycle, W-cycle, and full multigrid]{Three multigrid schedules, levels stacked
vertically (finest at top). \emph{Left:} the V-cycle makes one recursive call per level---one
trip down, one trip up. \emph{Middle:} the W-cycle makes two recursive calls per level,
visiting each coarse level twice for a stronger (but costlier) coarse correction. \emph{Right:}
full multigrid (FMG) starts on the coarsest grid and climbs, prolonging each level's solution
as the initial guess for a V-cycle on the next finer level---reaching discretization accuracy
in one upward sweep. The V-cycle is the workhorse; the others trade more coarse work for a
better correction.}
\label{fig:cycleshapes}
\end{figure}

\subsection{Auxiliary-space multigrid for $\Hcurl$}

The smoother-and-coarsen recipe assumes the smoother can damp all the oscillatory error---and
for the scalar Laplacian on $\Hone$ it can. But the curl--curl operator of every magnetic and
wave analysis (\cref{ch:functionspaces}) lives on $\Hcurl$, and there a plain smoother has a
blind spot: the large null space of the curl (the \emph{gradient fields}, $\Curl\Grad\phi =
0$) carries error that point smoothing does not see. A pure geometric V-cycle stalls on those
gradient modes. The cure is \emph{auxiliary-space} multigrid: alongside the geometric
hierarchy, build an auxiliary $\Hone$ space for the gradient part and smooth there too. The
per-level smoother becomes a \emph{Hiptmair smoother} (\cref{ch:smoothers}): a point smoother
on $\Hcurl$ to handle the non-gradient error, \emph{plus} a smoother on the auxiliary
gradient space to handle the kernel error the first one misses. Palace's AMS preconditioner
builds this auxiliary hierarchy automatically from the discrete gradient operator; the V-cycle
structure is unchanged---only the per-level smoother $\mat{B}_\ell$ is now the composite
Hiptmair one. The same shape, a better smoother, because the operator has a structure the
naive smoother cannot see.

\subsection{$h$- versus $p$-multigrid}

The hierarchy of \cref{fig:hierarchy} coarsened the \emph{mesh} ($h$-multigrid). One can
instead coarsen the \emph{polynomial order} on a fixed mesh ($p$-multigrid): the finest level
uses order-$p$ elements, the next order-$p/2$, down to order $1$ (or the family's floor from
\cref{ch:functionspaces}). The two are complementary and often combined ($hp$-multigrid):
coarsen the order down to linears, then coarsen the mesh. Palace's default schedule is
$p$-multigrid with a logarithmic order drop---halving the order each level---so a finest order
$p=8$ space reaches the floor in a handful of levels rather than the seven a linear
$p\mapsto p-1$ drop would need. The choice of $h$ versus $p$ changes \emph{what} the
prolongations interpolate (finer mesh, or higher order); the V-cycle recursion is identical.

\subsection{Multigrid as a preconditioner}

Here is the crucial deployment fact. We have described the V-cycle as a stand-alone iterative
solver---repeat V-cycles until the residual is small. In practice Palace almost never uses it
that way. Instead it uses \emph{one} V-cycle as a \emph{preconditioner} inside a Krylov method
(\cref{ch:cg,ch:krylov}). Recall from \cref{ch:precond} that a preconditioner is an
approximate inverse $\mat{M}^{-1}\approx\mat{A}^{-1}$ that the Krylov method applies once per
iteration to reshape the spectrum. \emph{One V-cycle is exactly such an $\mat{M}^{-1}$ apply}:
feed it a residual $\vr$ (as the right-hand side $\vb$, with zero initial guess) and it returns
an approximate correction $\mat{M}^{-1}\vr$---a good approximation to $\mat{A}^{-1}\vr$,
computed in $O(N)$ work. \Cref{fig:mgprecond} shows the arrangement: the Krylov solver drives
the outer loop, and each of its iterations calls one V-cycle to precondition.

\begin{figure}[t]
\centering
\begin{tikzpicture}[every node/.style={font=\footnotesize}]
  % outer Krylov loop with an inner V-cycle as the M^{-1} apply
  \node[stage, minimum width=30mm, minimum height=22mm] (krylov) at (0,0)
    {\textbf{Krylov solver}\\[2pt]\footnotesize (CG / GMRES)\\\footnotesize outer iteration};
  \node[extern, minimum width=30mm, minimum height=14mm, right=34mm of krylov] (vcyc)
    {\textbf{one V-cycle}\\[1pt]\footnotesize $= \mat{M}^{-1}$ apply\\\footnotesize $O(N)$ work};
  % residual goes in, correction comes back
  \draw[flow] (krylov.east|-0,0.35) -- node[above,font=\scriptsize]{residual $\vr$}
    (vcyc.west|-0,0.35);
  \draw[backflow] (vcyc.west|-0,-0.35) -- node[below,font=\scriptsize]{correction $\mat{M}^{-1}\vr$}
    (krylov.east|-0,-0.35);
  % the V-shape hint inside the extern box
  \node[anchor=north, simGray, font=\scriptsize\itshape] at (vcyc.south) {(the recursion of Fig.~\ref{fig:vcycle})};
  \node[anchor=south, simGray, font=\scriptsize\itshape] at (krylov.north)
    {once per iteration};
\end{tikzpicture}
\caption[Multigrid as a preconditioner]{How Palace actually deploys multigrid. A Krylov solver
(\cref{ch:cg,ch:krylov}) runs the outer iteration; once per iteration it hands its current
residual $\vr$ to \emph{one} V-cycle, which returns an approximate correction $\mat{M}^{-1}\vr$
in $O(N)$ work---the V-cycle \emph{is} the preconditioner's apply. Because the V-cycle makes
$\mat{M}^{-1}\mat{A}$ have an $h$-independent, tightly-clustered spectrum, the outer Krylov
iteration count is small and flat (\cref{fig:meshindep}). The combination---multigrid-preconditioned
CG or GMRES---is the standard optimal solver for the elliptic systems of \cref{part:modalities}.}
\label{fig:mgprecond}
\end{figure}

Why precondition rather than iterate V-cycles directly? Because the Krylov method is more
robust to imperfections in the V-cycle. If the smoother or coarsening is slightly off and a
few error modes converge slowly, a stand-alone V-cycle iteration inherits that slowness, but
the Krylov method \emph{accelerates} those stragglers with its optimal polynomial (the
machinery of \cref{ch:cg}). One V-cycle per Krylov iteration gives the best of both: the
V-cycle supplies the $h$-independent spectral clustering, and the Krylov method mops up
whatever the V-cycle handled imperfectly. This is the arrangement behind essentially every
elliptic solve in the book.

\section{Pitfalls: what breaks optimality}

Multigrid's optimality is not automatic; it rests on the two mechanisms---smoothing and
coarse-grid approximation---\emph{each} doing its job. Break either and the $h$-independent
convergence is lost, often silently, the method degrading into something no better than a
plain smoother.

\begin{pitfall}
\textbf{A bad smoother destroys optimality.} The whole scheme rests on the smoother
annihilating the error the coarse grid \emph{cannot} represent---the oscillatory modes. If the
smoother does not damp \emph{all} the oscillatory modes (e.g.\ an undamped Jacobi iteration,
which \emph{amplifies} the highest mode; or a point smoother on $\Hcurl$, blind to the gradient
null space of \cref{sec:variants}), those modes are left for a coarse grid that cannot see them,
and they never go away. The convergence factor $\rho$ creeps toward $1$ and the mesh-independence
evaporates. \emph{Always verify the smoother damps the full oscillatory band}---this is why
$\Hcurl$ problems need the Hiptmair smoother, not plain Jacobi.
\end{pitfall}

\begin{pitfall}
\textbf{A bad coarse grid destroys optimality too.} The coarse grid must \emph{capture the
smooth modes}---the error the smoother leaves behind. If the prolongation $\mat{P}$ cannot
represent the smooth error well (too aggressive a coarsening, a coarse space that omits a
near-null-space mode like the constant or a rigid-body mode), the coarse-grid correction is
inaccurate, the smooth error is not removed, and again $\rho\to 1$. The Galerkin construction
$\mat{A}_c=\mat{P}^{\mathsf T}\mat{A}\mat{P}$ guarantees the \emph{best} correction in the
range of $\mat{P}$, but only the range of $\mat{P}$---if the smooth error is not in that range,
no coarse solve recovers it.
\end{pitfall}

\begin{pitfall}
\textbf{Anisotropy and jumps need special handling.} The complementary-frequencies argument of
\cref{sec:complementary} assumed the smoother smooths \emph{isotropically}---equally in every
direction. For a strongly anisotropic operator (e.g.\ a Laplacian stretched
$1000\times$ in one direction, common in thin-layer geometries) a point smoother smooths only
along the strongly-coupled direction and leaves the other rough; standard coarsening then
coarsens a direction the smoother never smoothed, and optimality fails. The fixes---line/plane
smoothers that smooth a whole coupled direction at once, or semi-coarsening that coarsens only
the smoothed direction---are essential for anisotropic or jumping-coefficient problems and are
where naive multigrid most often disappoints. The lesson: the smoother and the coarsening must
be \emph{matched} to the operator's structure, not chosen in isolation.
\end{pitfall}

\summarysection
A smoother (\cref{ch:smoothers}) cheaply removes the high-frequency part of the error and
stalls on the smooth part. Multigrid turns that stall into a method by exploiting the
\emph{complementary-frequencies principle}: smooth error on a fine grid looks oscillatory on a
coarser grid, where a smoother can remove it cheaply. The two-grid method makes this concrete:
presmooth on the fine grid, compute the residual, restrict it to the coarse grid with
$\mat{R}=\mat{P}^{\mathsf T}$, solve the Galerkin coarse problem $\mat{A}_c=\mat{P}^{\mathsf
T}\mat{A}\mat{P}$ for a correction, prolong the correction back with $\mat{P}$, and postsmooth.
The \emph{V-cycle} applies this idea recursively: each level's coarse solve is itself a V-cycle
on a still-coarser grid, until a tiny coarsest grid is solved directly---a descending stroke of
presmoothing-and-restricting and an ascending stroke of prolonging-and-postsmoothing, the V of
the hero diagram. Written as a pure recursive function, every leg is the same correction step
$\vy + \mat{B}(\vb-\mat{A}\vy)$. For finite elements the prolongations come free from the nested
function-space hierarchy of \cref{ch:functionspaces}. The V-cycle is \emph{optimal}: one cycle
costs $O(N)$ (the per-level costs form a geometric series summing to a small constant times $N$),
and a fixed, $h$-independent number of cycles reaches any tolerance---in sharp contrast to
unpreconditioned CG, whose iteration count grows like $1/h$. Variants trade coarse work for
correction strength (W-cycle, full multigrid) or adapt the smoother to the operator (Hiptmair
relaxation for the $\Hcurl$ gradient null space; $p$- versus $h$-coarsening). In practice one
V-cycle is used as the $\mat{M}^{-1}$ apply inside a Krylov solver---the optimal preconditioner
of \cref{ch:precond}. Optimality is fragile: a smoother that misses some oscillatory modes, a
coarse grid that misses some smooth modes, or an anisotropic operator that defeats isotropic
smoothing each silently destroys the mesh-independence.

\furtherreading
The auxiliary-space and Hiptmair smoother machinery for $\Hcurl$ multigrid is developed in
Hiptmair's foundational paper~\citep{hiptmair1998}, the source of the gradient-space smoother
that makes curl--curl multigrid work; it is the natural sequel to \cref{sec:variants}. For the
finite-element and variational framing of multigrid---nested spaces, Galerkin coarsening, the
approximation-and-smoothing convergence theory sketched in \cref{sec:optimal}---Brenner and
Scott~\citep{brennerscott2008} give the rigorous account at the level of this book. A modern,
implementation-oriented survey of multigrid and $p$-multigrid for high-order finite elements,
close to Palace's actual algorithms, is Phillips and Fischer~\citep{phillipsfischer2022}. The
classical introduction to the complementary-frequencies argument, with the Fourier (local-mode)
analysis of smoothing factors only sketched here, remains the multigrid tutorial literature
those references cite in turn.

\exercisessection
\begin{enumerate}[nosep]
\item State the complementary-frequencies principle in one sentence, and explain why it would
  fail if the smoother damped the \emph{low} frequencies rather than the high ones. What kind of
  method would you need to pair with such a smoother instead of a coarse grid?
\item For the 1-D linear-interpolation prolongation of \cref{sec:twogrid}, write out the
  $7\times 3$ matrix $\mat{P}$ mapping a coarse grid of $3$ interior points to a fine grid of $7$
  interior points (spacing $2h$ to $h$). Then write $\mat{R}=\mat{P}^{\mathsf T}$ and confirm it
  is the full-weighting $\tfrac14[1,2,1]$ restriction stencil up to scaling.
\item Using the $\mat{P}$ from the previous exercise and the fine tridiagonal Laplacian
  $\mat{A}=h^{-2}\mathrm{tridiag}(-1,2,-1)$, compute the Galerkin coarse operator
  $\mat{A}_c=\mat{P}^{\mathsf T}\mat{A}\mat{P}$ and verify it is (up to scaling) the $3\times 3$
  tridiagonal Laplacian at spacing $2h$. What does this say about Galerkin coarsening for this
  problem?
\item Prove \cref{lem:galerkin} for yourself without looking: show that if $\mat{A}$ is SPD and
  $\mat{P}$ has full column rank, then $\mat{A}_c=\mat{P}^{\mathsf T}\mat{A}\mat{P}$ is SPD. Which
  hypothesis fails, and what goes wrong, if $\mat{P}$ has a nontrivial null space?
\item A 3-D problem has $N = 10^6$ unknowns on the finest grid and $5$ levels. Using the
  geometric-series estimate of \cref{wex:cost}, compute the total unknown count summed over all
  levels and express it as a multiple of $N$. How much does the entire coarse hierarchy add to the
  finest-level cost?
\item Unpreconditioned CG on a 3-D Poisson problem needs $\approx\tfrac12\sqrt{\kappa}$ iterations
  with $\kappa\approx (1/h)^2$. If multigrid-preconditioned CG needs a flat $\approx 10$ iterations
  regardless of $h$, tabulate the iteration count of each method for meshes
  $h=\tfrac1{16},\tfrac1{32},\tfrac1{64},\tfrac1{128}$. By what factor does plain CG's count grow
  per refinement, and multigrid's?
\item Explain, in terms of the V-cycle recursion of \cref{lst:vcycle}, the difference between a
  V-cycle and a W-cycle. Modify the pseudocode of \cref{lst:vcycle} to perform a W-cycle (two
  recursive calls per level). Why does the W-cycle cost more, and when is the extra cost worth it?
\item Why does a point Jacobi smoother stall on the curl--curl operator's gradient null space, and
  how does the Hiptmair smoother of \cref{sec:variants} fix it? Relate your answer to the
  spurious-mode discussion of \cref{ch:functionspaces}: both are consequences of the structure of
  $\Hcurl$.
\end{enumerate}
